\begin{tikzpicture}[x=8mm]
  % sposób 1: posortowane litery
  \foreach \c [count=\i from 0] in {l,a,m,p,a} \node[komorka] (x\i) at (\i, 0) {\c};
  \foreach \c [count=\i from 0] in {p,a,l,m,a} \node[komorka] (y\i) at (\i, -1.3) {\c};
  \node[kod, anchor=east] at (-0.7, 0) {"lampa"};
  \node[kod, anchor=east] at (-0.7, -1.3) {"palma"};
  \draw[strzalka] (4.7, 0) -- node[above, font=\scriptsize\ttfamily] {sorted} (6.6, 0);
  \draw[strzalka] (4.7, -1.3) -- node[above, font=\scriptsize\ttfamily] {sorted} (6.6, -1.3);
  \foreach \c [count=\i from 0] in {a,a,l,m,p} {
    \node[komorka ok] (sx\i) at (7.3+\i, 0) {\c};
    \node[komorka ok] (sy\i) at (7.3+\i, -1.3) {\c};
    \node[font=\scriptsize, text=zielen] at (7.3+\i, -0.65) {=};
  }
  \node[podpis, anchor=north] at (9.3, -1.9) {te same listy → anagramy};
  % sposób 2: liczniki liter
  \begin{scope}[xshift=13.6*8mm]
    \foreach \l/\h [count=\i from 0] in {a/2,l/1,m/1,p/1} {
      \fill[akcentjasny] (\i*0.8-0.25, -1.6) rectangle (\i*0.8+0.25, -1.6+\h*0.7);
      \draw[akcent] (\i*0.8-0.25, -1.6) rectangle (\i*0.8+0.25, -1.6+\h*0.7);
      \node[kod, font=\ttfamily\small] at (\i*0.8, -1.85) {\l};
      \node[font=\scriptsize] at (\i*0.8, -1.45+\h*0.7) {\h};
    }
    \draw[draw=szary] (-0.45, -1.6) -- (2.85, -1.6);
    \node[podpis, anchor=north, text width=36mm] at (1.2, -2.05) {w \textbf{obu} słowach każda\\ litera występuje\\ tyle samo razy};
  \end{scope}
\end{tikzpicture}
